% !TeX root = ../navier-stokes-manuscript.tex
\subsection{Why the finite pieces are rectangles}\label{sub:BodyFiniteAdjacentRectangles}

Keep the maximal classical solution generated by the fixed pair \((u_0,\nu)\), and write its lifespan as \([0,T_*)\).
On the branch where that lifespan is finite, set
\[
  I_*:=(0,T_*).
\]
A spatial norm becomes useful at the endpoint only if it has a time trace.
That trace comes from placing a response beside its time derivative.
Let \(A=(1-\Delta)^{1/2}\).  For a smooth Hilbert-valued function \(v\) and \(0\leq a<b\leq T\),
\begin{equation}\label{eq:BodyAdjacentTraceIdentity}
  \norm{v(b)}_{H^M}^2-
  \norm{v(a)}_{H^M}^2
  =2\operatorname{Re}
  \int_a^b
  \pair{A^{M-1}v_t(t)}{A^{M+1}v(t)}_{L^2}\,dt.
\end{equation}
Cauchy--Schwarz places the change of the middle norm between two adjacent space--time rows:
\[
  \left|
  \norm{v(b)}_{H^M}^2-
  \norm{v(a)}_{H^M}^2
  \right|
  \leq
  2\norm{v_t}_{L^2(a,b;H^{M-1})}
   \norm{v}_{L^2(a,b;H^{M+1})}.
\]
Consequently,
\begin{equation}\label{eq:BodyAdjacentTracePair}
  v\in L^2(0,T;H^{M+1}),
  \qquad
  v_t\in L^2(0,T;H^{M-1})
\end{equation}
gives
\begin{equation}\label{eq:BodyAdjacentTraceConclusion}
  v\in C([0,T];H^M).
\end{equation}
The proof in \Cref{lem:AdjacentHilbertScaleTrace} uses time mollification and
Fourier truncation to justify this identity; weak \(H^M\) continuity together
with continuity of the norm gives the strong representative.

At \(M=2\), the mechanism reads
\begin{equation}\label{eq:FirstConcreteAdjacentPair}
  u\in L^2(0,T;H^3),
  \qquad
  u_t\in L^2(0,T;H^1)
  \quad\Longrightarrow\quad
  u\in C([0,T];H^2).
\end{equation}
The velocity supplies one derivative above the desired trace, its time response supplies one derivative below, and their pairing reaches the middle space.

Apply the same trace to every \(V_n\), with \(v_t=V_{n+1}\).
To bring the first \(N+1\) responses into \(H^M\), collect exactly the adjacent rows they require:
\begin{equation}\label{eq:BodyRectangleNorm}
\begin{aligned}
  \mathfrak R_{N,M}(u;T)
  :={}&
  \sum_{n=0}^{N}
  \norm{V_n}_{L^2(0,T;H^{M+1})}^2
  \\
  &+
  \sum_{n=0}^{N}
  \norm{V_{n+1}}_{L^2(0,T;H^{M-1})}^2.
\end{aligned}
\end{equation}
When \(M=0\), the lower edge is read in \(H^{-1}\).
The temporal depth \(N\) and spatial height \(M\) remain independent choices.

\Needspace{12\baselineskip}
For example, \((N,M)=(1,2)\) contains two trace pairs sharing the row \(V_1\):
\begin{center}
\captionsetup{hypcap=false}
\captionof{table}{The \((1,2)\) finite rectangle.  The row \(V_1\) is shared: it is the time derivative used to trace \(V_0\), and it is also the next response whose own trace is paired with \(V_2\).}
\label{tab:FirstAdjacentRectangle}
\renewcommand{\arraystretch}{1.35}
\begin{tabular}{@{}c c c@{}}
\toprule
temporal row & one step above & one step below \\
\midrule
\(n=0\) & \(V_0=u\in L^2H^3\) & \(V_1=u_t\in L^2H^1\) \\
\(n=1\) & \(V_1=u_t\in L^2H^3\) & \(V_2=u_{tt}\in L^2H^1\) \\
\bottomrule
\end{tabular}
\end{center}

Write \(\cR_{N,M}[u_0;T]\) for this finite pressure-complete block: the velocity coordinates measured by \eqref{eq:BodyRectangleNorm} together with the corresponding \(Q_n\)-coordinates reconstructed by \eqref{eq:BodyMixedPressureRecurrence}.
Every entry remains a derivative of the datum-generated velocity and pressure fields.

For a strict cutoff \(T<T_*\), classical smoothness makes every fixed block finite:
\begin{equation}\label{eq:StrictSubintervalRectangleBound}
  \mathfrak R_{N,M}(u;T)<\infty
  \qquad\text{for finite }N,M.
\end{equation}
But it supplies only the left-hand quantifier:
\[
  \forall\,T<T_*:\quad \mathfrak R_{N,M}(u;T)<\infty
  \quad\not\Longrightarrow\quad
  \sup_{0<T<T_*}\mathfrak R_{N,M}(u;T)<\infty.
\]
The missing estimate reverses these quantifiers.
After \(N\) and \(M\) are fixed, one bound must survive every cutoff approaching \(T_*\).
